\section{Stability tricks}

本节处理 softmax 和 logits 的稳定性：z-loss、QK norm、logit soft-capping。


\begin{warningbox}{稳定性技巧的共同风险}
z-loss、QK norm、logit soft-capping 都在限制数值尺度。它们能减少训练爆炸或 softmax 饱和，但也可能改变模型表达范围。稳定性补丁必须和 loss 曲线、梯度 norm、下游指标一起看。
\end{warningbox}


超参数 recap 之后，Slide 52 用一条平稳曲线和一条爆炸曲线界定本节目标：稳定性技巧不是为了让正常训练再快一点，而是避免少数 step 的数值失控毁掉昂贵 run。大模型、长上下文和低精度训练会放大 logits、attention scores 与 gradient spikes，因此必须区分“优化不佳”和“数值范围已经失控”。

\begin{figure}[H]
\centering
\includegraphics[width=0.88\textwidth]{slide-051.jpg}
\caption{Slide 52：稳定训练与爆炸训练曲线；后续技巧都在控制数值尺度。}
\figsource{CS336 Spring 2026 Lecture 3 官方幻灯片。}
\end{figure}

\begin{knowledgebox}{稳定性排查先看哪几条时间序列}
至少同时记录 training loss、gradient norm、logit RMS/max、attention score RMS/max、activation RMS 和 learning rate。若 loss 突然上升前 logits 或 QK score 已持续增大，z-loss、QK norm 或 soft-capping 才是针对性工具；若首先出现通信错误、NaN kernel 或数据异常，则不应把问题归因于 architecture。
\end{knowledgebox}


\begin{figure}[H]
\centering
\includegraphics[width=0.88\textwidth]{slide-052.jpg}
\caption{Slide 53: Softmax issues}
\figsource{CS336 Spring 2026 Lecture 3 官方幻灯片。}
\end{figure}

\noindent\textbf{展开说明：}本页指出 softmax 因指数和除法容易数值不稳定，是 output 和 attention 两处的共同风险。


\begin{figure}[H]
\centering
\includegraphics[width=0.88\textwidth]{slide-053.jpg}
\caption{Slide 54: Z-loss}
\figsource{CS336 Spring 2026 Lecture 3 官方幻灯片。}
\end{figure}

\noindent\textbf{展开说明：}本页定义 z-loss：约束 log-sum-exp normalizer，防止 logits 尺度失控。

\begin{knowledgebox}{读图：Slide 54 应该怎么看}
这页不是装饰性截图，而是本节论点的证据页。先看图中模型/论文/曲线或表格的比较对象，再看它们支持的趋势：哪些设计已经形成共识，哪些只是局部实验结果，哪些主要由运行时或推理成本推动。读这类页时要区分 quality evidence、runtime evidence 和 stability evidence，不能把一种证据直接替换成另一种。
\end{knowledgebox}


\begin{figure}[H]
\centering
\includegraphics[width=0.88\textwidth]{slide-054.jpg}
\caption{Slide 55: QK norm}
\figsource{CS336 Spring 2026 Lecture 3 官方幻灯片。}
\end{figure}

\noindent\textbf{展开说明：}本页讲 QK norm：对 query/key 做 norm 控制 attention logits，减少 softmax 极端值。

\begin{knowledgebox}{读图：Slide 55 应该怎么看}
这页不是装饰性截图，而是本节论点的证据页。先看图中模型/论文/曲线或表格的比较对象，再看它们支持的趋势：哪些设计已经形成共识，哪些只是局部实验结果，哪些主要由运行时或推理成本推动。读这类页时要区分 quality evidence、runtime evidence 和 stability evidence，不能把一种证据直接替换成另一种。
\end{knowledgebox}


\begin{figure}[H]
\centering
\includegraphics[width=0.88\textwidth]{slide-055.jpg}
\caption{Slide 56: Logit soft-capping}
\figsource{CS336 Spring 2026 Lecture 3 官方幻灯片。}
\end{figure}

\noindent\textbf{展开说明：}本页展示用 tanh 对 logits 做 soft cap，限制极端 logits。

\subsection{本章小结}
本节的关键不是记住所有模型名，而是把公开模型的设计选择转化成可迁移判断：共识默认值可以作为起点，例外需要知道原因，涉及系统和推理成本的选择必须结合资源账本讨论。

